% ============================================================================
% ABLATION: mass-function-informed population vs free-Gaussian population.
% Experiment: obs_z1_lambda5_high_rm_scatter.376 (rm_scatter = 0.7).
% Numbers from notebooks/hier_massfn_informed.py -> massfn_informed.pkl (run on main thread).
% Figure: massfn_informed_compare.png ; label \ref{fig:massfn} wired by main thread.
% RESULT: informed prior does NOT help -- a null/cautionary result.
% ============================================================================

As an ablation we relax the hierarchical model's deliberate model-independence and ask what is gained,
and what is given up, by informing the mass population with external physics. We replace the free
Gaussian mass population $\log_{10}M_j \sim \mathcal{N}(\mu_M, \sigma_M)$ with a population whose
per-cluster masses carry an informative prior equal to the in-bin mass distribution predicted by the
\citet{tinker08} halo mass function multiplied by the \citet{mcclintock18} richness--mass
selection window for $30 < \lambda < 45$ (the same mass-function and selection terms as our reference
informed MCMC, with log-normal richness scatter $\sigma_{\ln\lambda}=\sigma_{\log M}\ln 10/F$,
$F=1.356$). The concentration sector ($c_0,\beta,\sigma_c$) and the profile likelihood are held fixed,
so any change is attributable solely to the mass prior. We run this on the high $\lambda$--$M$ scatter
experiment, where the true in-bin mass distribution is most strongly non-Gaussian (true excess
kurtosis $-0.51$) and where the free model recovers the slope $\beta$ worst.

The ablation returns a clear \emph{null} result, which is itself informative. On every recoverable
quantity the informed population matches the free Gaussian rather than improving on it: the recovered
mass spread is $\sigma_M=0.49$ (informed) versus $0.52$ (free), against a true $0.535$---the informed
prior is in fact marginally \emph{worse}---and both already reproduce the non-Gaussian, flat-topped
shape of the true distribution (recovered excess kurtosis $-0.53$ informed, $-0.48$ free, versus true
$-0.51$; Figure~\ref{fig:massfn}). The free hierarchical model, despite assuming a Gaussian
\emph{prior} on the population, recovers a non-Gaussian \emph{posterior} population because the
per-cluster latents are free to populate the tails---so the Gaussian-population assumption is far less
restrictive in practice than it appears. For the weakly identified slope, the informed prior does not
help and slightly hurts: $\beta=-1.04\pm0.18$ (informed) versus $-1.28\pm0.17$ (free) against true
$-1.33$, with no reduction in posterior width. This is consistent with our finding that $\beta$ is
limited by the mass lever arm (Section~\ref{sec:beta}, Figure~\ref{fig:beta_leverarm}), not by the
mass prior: pinning the mass distribution to the mass function cannot supply the joint $(M,c)$
information that the measurement noise removes. The practical conclusion is that informing the
population with the mass function and richness--mass relation buys essentially nothing here while
reintroducing exactly the model-dependence---the assumed mass function and selection function---that
the free hierarchical approach was designed to avoid. We therefore retain the free population as our
default and report this ablation as a deliberate contrast.

\begin{figure}[t]
\centering
\includegraphics[width=\columnwidth]{figures/massfn_informed_compare.png}
\caption{Mass-function-informed vs free-Gaussian population, high $\lambda$--$M$ scatter experiment.
\emph{Left:} recovered in-bin mass distribution for the free-Gaussian (blue) and mass-function-informed
(green) hierarchical models against the true drawn population (grey); both reproduce the non-Gaussian,
flat-topped truth, and both differ markedly from the raw mass-function$\times$selection prior (red
dashed)---the data, not the prior, drive the result. \emph{Right:} the $M$--$c$ slope posterior: the
free model (blue) sits on the true value (dashed), while the informed prior pulls the informed model
(green) \emph{away} from it.\label{fig:massfn}}
\end{figure}
